% Exploratory Results shells, written without reading any sealed output.
% Governed by pipeline/review/x3-x7-reading-restrictions-2026-09-29.md.
% Not connected to main.tex. Each numbered branch below is a prose option.
% Copy the status paragraphs, applicable branches and table inputs into the
% supplement after release. Omit zero-count outcome branches. Every displayed
% number and machine state comes from v4_exploratory_export; fill from its
% catalog, optionally merged with the confirmatory catalog. No branch chooses
% a cell, changes a threshold, or promotes exploratory evidence.

% STATUS: always; incomplete counts refer to the declared operational cut.
\subsection*{Exploratory analyses of fitting, rejection, coverage and features}\label{supp:x3-x7}
\paragraph{Status and scope.}
These analyses used declared exploratory seeds and leave the frozen confirmatory design unchanged.
X3 recorded \resultnum{v4x.x3.recorded}{[n]} \resultplural{v4x.x3.recorded}{cell}{cells}, with an incomplete count of \resultnum{v4x.x3.incomplete}{[n]} \resultplural{v4x.x3.incomplete}{cell}{cells}, \resultnum{v4x.x3.errors}{[n]} execution \resultplural{v4x.x3.errors}{error}{errors} and \resultnum{v4x.x3.gaps}{[n]} fit \resultplural{v4x.x3.gaps}{gap}{gaps}.
X4a recorded \resultnum{v4x.x4a.recorded}{[n]} \resultplural{v4x.x4a.recorded}{cell}{cells}, with an incomplete count of \resultnum{v4x.x4a.incomplete}{[n]}. X4b recorded \resultnum{v4x.x4b.recorded}{[n]} \resultplural{v4x.x4b.recorded}{cell}{cells}, with an incomplete count of \resultnum{v4x.x4b.incomplete}{[n]}.
X4a had \resultnum{v4x.x4a.errors}{[n]} execution \resultplural{v4x.x4a.errors}{error}{errors} and \resultnum{v4x.x4a.gaps}{[n]} fit \resultplural{v4x.x4a.gaps}{gap}{gaps}; X4b had \resultnum{v4x.x4b.errors}{[n]} execution \resultplural{v4x.x4b.errors}{error}{errors} and \resultnum{v4x.x4b.gaps}{[n]} fit \resultplural{v4x.x4b.gaps}{gap}{gaps}.
The cellwise contrasts have no multiplicity correction or whole-curve reading (Table~\ref{tab:v4x-status}).

% X3-DESIGN: always with X3; the design and the common X3 caveat, stated once.
\paragraph{X3: design and common limits.}
X3 lets null B choose its own response signal share. In the free-share arm, B's frozen latent law passes through the observed channel at a fitted share $r_B$, matched together with B's two persistence parameters to the lag-one to lag-three moments. Because lag three becomes a matched moment, both X3 contrasts train the LDA on the lag-four and lag-five features only, unlike the main discriminator, which uses lags three to five. The known-share contrast uses the frozen two-parameter B at the cell's own $r$ and the free-share contrast uses the fitted $r_B$, both on the same social panels; each cell reports the free-minus-known AUC difference with a paired bootstrap interval, read against the declared margin of $\pm\resultnum{v4x.x3.summary.margin}{[value]}$.
These intervals are conditional on the accepted fits and the trained discriminators and omit the uncertainty of the three-parameter matching; the contrasts do not estimate the frozen boundary.

% X3-SUPPORT: count.supports-objection > 0.
\paragraph{X3: reduced separation.}
Allowing B its own response signal share reduced separability for the specified L4--L5 LDA in \resultnum{v4x.x3.count.supports-objection}{[n]} \resultplural{v4x.x3.count.supports-objection}{cell}{cells}. \resultplural{v4x.x3.count.supports-objection}{This cell had}{These cells had} a negative paired AUC-difference interval (reported label \texttt{supports-objection}; Table~\ref{tab:v4x-x3}).

% X3-EQUIVALENCE: count.against-objection > 0.
\paragraph{X3: local equivalence.}
In \resultnum{v4x.x3.count.against-objection}{[n]} \resultplural{v4x.x3.count.against-objection}{cell}{cells}, the paired interval lay within the declared margin (\texttt{against-objection}).
This is local equivalence for the specified contrast; it does not establish equivalence of mechanisms or of their full laws.

% X3-INCONCLUSIVE: count.inconclusive > 0.
\paragraph{X3: inconclusive contrast.}
The cellwise contrast was inconclusive in \resultnum{v4x.x3.count.inconclusive}{[n]} \resultplural{v4x.x3.count.inconclusive}{cell}{cells} (\texttt{inconclusive}); \resultplural{v4x.x3.count.inconclusive}{the interval in}{the intervals in} Table~\ref{tab:v4x-x3} \resultplural{v4x.x3.count.inconclusive}{retains}{retain} the uncertainty under the declared conditional analysis.

% X3-ERROR: count.error > 0.
\paragraph{X3: execution errors.}
There \resultplural{v4x.x3.count.error}{was}{were} \resultnum{v4x.x3.count.error}{[n]} recorded execution \resultplural{v4x.x3.count.error}{error}{errors} (\texttt{error}).
\resultplural{v4x.x3.count.error}{The affected cell provides}{The affected cells provide} no AUC contrast.

% X3-GAP: count.gap > 0.
\paragraph{X3: known-fit gaps.}
There \resultplural{v4x.x3.count.gap}{was}{were} \resultnum{v4x.x3.count.gap}{[n]} known-fit \resultplural{v4x.x3.count.gap}{gap}{gaps} (\texttt{gap}).
Both contrasts remain unavailable in those cells, with the known-fit attempts retained.

% X3-FREE-GAP: count.free-gap > 0.
\paragraph{X3: free-fit gaps.}
There \resultplural{v4x.x3.count.free-gap}{was}{were} \resultnum{v4x.x3.count.free-gap}{[n]} free-fit \resultplural{v4x.x3.count.free-gap}{gap}{gaps} (\texttt{free-gap}).
Table~\ref{tab:v4x-x3-attempts} reports each recorded attempt, its matching size and the last free-fit status. The accompanying values catalog reports residuals and evaluations.
Search budget exhaustion, a numerical search without an accepted fit and a candidate rejected by validation remain distinct outcomes.
A family-limit or no-root search label alone does not demonstrate that the free-share family cannot reproduce the target moments.

% X3-LOST: count.lost_crossing > 0.
\paragraph{X3: paired-cell flag.}
In \resultnum{v4x.x3.count.lost_crossing}{[n]} \resultplural{v4x.x3.count.lost_crossing}{cell}{cells}, the known-share point AUC was at least \resultnum{v4.protocol.target}{[value]} while the free-share interval upper endpoint was below that target (\texttt{lost\_crossing}).
This is the declared paired-cell flag, without a frozen-band crossing or a confirmatory boundary claim.

% X4-STATUS: always.
\paragraph{X4: conditional rejection references.}
The oracle direction knows the alternative.
Its critical value was estimated from \resultnum{v4x.x4.sizes.check_prose}{[n]} null \resultplural{v4x.x4.sizes.check_prose}{score}{scores}. The fitted-panel procedure used \resultnum{v4x.x4.sizes.analyst}{[n]} simulated \resultplural{v4x.x4.sizes.analyst}{panel}{panels} per fitting evaluation and \resultnum{v4x.x4.sizes.nulls}{[n]} fresh null \resultplural{v4x.x4.sizes.nulls}{panel}{panels} for its bootstrap.
Tables~\ref{tab:v4x-x4a} and \ref{tab:v4x-x4b} report rejection references, intervals and panel denominators. Table~\ref{tab:v4x-x4-failures} separately reports failures counted as non-rejections and the success-conditioned reference.
The bootstrap p-value formula establishes neither exact validity over the composite null nor an analyst's attainable power.
X4b and its zero-strength calibration reference are declared only at the three response signal shares of Table~\ref{tab:v4x-x4b}; values outside this domain are not missing in-domain estimates.

% X4-NORMAL: each positive count is a separate cellwise branch.
\paragraph{X4: compatible normal reference.}
The normal-theory reference lay inside the oracle rejection interval in \resultnum{v4x.x4a.count.describes}{[n]} \resultplural{v4x.x4a.count.describes}{cell}{cells} within the declared reference-AUC range of \resultnum{v4x.x4.summary.normal_theory_in_auc_range.range.i000}{[lower]} to \resultnum{v4x.x4.summary.normal_theory_in_auc_range.range.i001}{[upper]} (\texttt{describes}; Table~\ref{tab:v4x-x4-normal}).
This records interval compatibility in those cells.

\paragraph{X4: higher normal reference.}
The normal-theory reference lay above the oracle rejection interval in \resultnum{v4x.x4a.count.overstates}{[n]} \resultplural{v4x.x4a.count.overstates}{cell}{cells} within the declared reference-AUC range of \resultnum{v4x.x4.summary.normal_theory_in_auc_range.range.i000}{[lower]} to \resultnum{v4x.x4.summary.normal_theory_in_auc_range.range.i001}{[upper]} (\texttt{overstates}).
The reference was evaluated at each cell's own AUC.

\paragraph{X4: lower normal reference.}
The normal-theory reference lay below the oracle rejection interval in \resultnum{v4x.x4a.count.understates}{[n]} \resultplural{v4x.x4a.count.understates}{cell}{cells} within the declared reference-AUC range of \resultnum{v4x.x4.summary.normal_theory_in_auc_range.range.i000}{[lower]} to \resultnum{v4x.x4.summary.normal_theory_in_auc_range.range.i001}{[upper]} (\texttt{understates}).
These cellwise comparisons do not validate the normal-theory translation.

% X4-PAIRED: each positive count is a separate branch.
\paragraph{X4: negative paired difference.}
The paired difference between the conditional oracle rejection procedures had a negative interval in \resultnum{v4x.x4b.count.removes-power}{[n]} \resultplural{v4x.x4b.count.removes-power}{cell}{cells} (\texttt{removes-power}; Table~\ref{tab:v4x-x4-paired}).
It combines fitted-panel estimation, simulation precision and failure treatment and does not isolate the effect of fitting alone.

\paragraph{X4: local equivalence of rejection references.}
The paired interval lay within the declared margin of $\pm\resultnum{v4x.x4.summary.single_panel_effect.margin}{[value]}$ in \resultnum{v4x.x4b.count.no-material-change}{[n]} \resultplural{v4x.x4b.count.no-material-change}{cell}{cells} (\texttt{no-material-change}).
This is local equivalence of the two conditional rejection procedures.

\paragraph{X4: inconclusive paired difference.}
The paired difference was inconclusive in \resultnum{v4x.x4b.count.inconclusive}{[n]} \resultplural{v4x.x4b.count.inconclusive}{cell}{cells} (\texttt{inconclusive}).
The size flags and original quantity labels remain beside the references; unavailable readings are retained as such.

% X4-SIZE: each positive count is a separate branch; counts are zero cells.
\paragraph{X4: zero-strength compatibility.}
Rejection was compatible with the nominal level of \resultnum{v4x.protocol.x4.level}{[value]} at \resultnum{v4x.x4a.count.calibrated}{[n]} X4a zero-strength \resultplural{v4x.x4a.count.calibrated}{cell}{cells} and \resultnum{v4x.x4b.count.calibrated}{[n]} X4b zero-strength \resultplural{v4x.x4b.count.calibrated}{cell}{cells} (\texttt{calibrated}; Table~\ref{tab:v4x-x4-calibration}).
Each interval and denominator applies to its specified cell, and compatibility there does not establish calibration over the composite null.

\paragraph{X4: above-nominal zero-strength reference.}
The above-nominal zero-strength reference counts under the declared criterion were \resultnum{v4x.x4a.count.liberal}{[n]} X4a and \resultnum{v4x.x4b.count.liberal}{[n]} X4b cells (\texttt{liberal}).
Positive-strength rejection references at those signal shares retain that local flag.

\paragraph{X4: below-nominal zero-strength reference.}
The below-nominal zero-strength reference counts under the declared criterion were \resultnum{v4x.x4a.count.conservative}{[n]} X4a and \resultnum{v4x.x4b.count.conservative}{[n]} X4b cells (\texttt{conservative}).
The original machine quantity is retained beside the conditional reference, without a universal bound on valid power.

% X6-STATUS: always for the inventory sentences. The joint coverage sentences
% require x6a.band.vs_d1 != null. Omit them if procedural truth is unavailable.
\paragraph{X6: coverage denominators.}
At the operational cut, X6a had recorded \resultnum{v4x.x6.x6a.recorded_prose}{[n]} of \resultnum{v4x.x6.x6a.units_prose}{[n]} \resultplural{v4x.x6.x6a.units_prose}{job}{jobs}, with an incomplete count of \resultnum{v4x.x6.x6a.incomplete}{[n]} \resultplural{v4x.x6.x6a.incomplete}{job}{jobs}. X6b had recorded \resultnum{v4x.x6.x6b.recorded}{[n]} of \resultnum{v4x.x6.x6b.units}{[n]} \resultplural{v4x.x6.x6b.units}{job}{jobs}, with an incomplete count of \resultnum{v4x.x6.x6b.incomplete}{[n]} \resultplural{v4x.x6.x6b.incomplete}{job}{jobs} (Table~\ref{tab:v4x-x6-inventory}).
Primary joint D1 coverage was \resultnum{v4x.x6.x6a.band.vs_d1.primary.covered}{[n]} of \resultnum{v4x.x6.x6a.band.vs_d1.primary.total}{[n]} complete gap-free \resultplural{v4x.x6.x6a.band.vs_d1.primary.total}{band}{bands} (Table~\ref{tab:v4x-x6}). Joint coverage requires coverage at \resultplural{v4x.protocol.grid_strengths}{the}{all} \resultnum{v4x.protocol.grid_strengths}{[n]} \resultplural{v4x.protocol.grid_strengths}{strength}{strengths}, including $a=\resultnum{v4x.protocol.zero}{[value]}$.
The separate gaps-as-misses reference was \resultnum{v4x.x6.x6a.band.vs_d1.gaps_as_misses.covered}{[n]} of \resultnum{v4x.x6.x6a.band.vs_d1.gaps_as_misses.total}{[n]} \resultplural{v4x.x6.x6a.band.vs_d1.gaps_as_misses.total}{repetition}{repetitions} with a curve.
There were \resultnum{v4x.x6.x6a.without_curve}{[n]} \resultplural{v4x.x6.x6a.without_curve}{repetition}{repetitions} without a curve and \resultnum{v4x.x6.x6a.repetitions.errors}{[n]} recorded repetition \resultplural{v4x.x6.x6a.repetitions.errors}{error}{errors}. These counts and incomplete jobs are reported outside the computed coverage denominators.
Table~\ref{tab:v4x-x6-denominators} states the conditioning for each cellwise reference: successful complete cell pipelines for X6b and complete gap-free units for the X6a primary readings.

% X6-READINGS: counts refer to tabulated D1 coverage references, never cells
% pooled as repetitions. Retain the particular Wilson interval from the table.
\paragraph{X6: no undercoverage detected.}
We did not detect undercoverage under the declared criterion in \resultnum{v4x.x6.x6a.count.adequate}{[n]} X6a and \resultnum{v4x.x6.x6b.count.adequate}{[n]} X6b coverage references against D1 (\texttt{adequate}).
Their Wilson upper endpoints were at least \resultnum{v4x.protocol.x6.nominal}{[value]}; Table~\ref{tab:v4x-x6} gives each covered count, denominator and interval.

\paragraph{X6: below-nominal reference.}
The counts were \resultnum{v4x.x6.x6a.count.subnominal}{[n]} X6a \resultplural{v4x.x6.x6a.count.subnominal}{reference}{references} and \resultnum{v4x.x6.x6b.count.subnominal}{[n]} X6b \resultplural{v4x.x6.x6b.count.subnominal}{reference}{references} with Wilson upper endpoints below \resultnum{v4x.protocol.x6.nominal}{[value]} but at least \resultnum{v4x.protocol.x6.failure}{[value]} against D1 (\texttt{subnominal}).
These intervals give evidence of undercoverage under that criterion in the stated denominators.

\paragraph{X6: failure criterion.}
The Wilson upper endpoint was below \resultnum{v4x.protocol.x6.failure}{[value]} in \resultnum{v4x.x6.x6a.count.failure}{[n]} X6a and \resultnum{v4x.x6.x6b.count.failure}{[n]} X6b coverage references against D1 (\texttt{failure}).
The count includes separately tabulated primary and gaps-as-misses rows against D1; it does not count distinct cells. Their covered counts and intervals retain their declared conditioning.

\paragraph{X6: no reading.}
No coverage reading was assigned to \resultnum{v4x.x6.x6a.count.no-reading}{[n]} X6a \resultplural{v4x.x6.x6a.count.no-reading}{reference}{references} and \resultnum{v4x.x6.x6b.count.no-reading}{[n]} X6b \resultplural{v4x.x6.x6b.count.no-reading}{reference}{references} because each fell below the minimum of \resultnum{v4x.x6.minimum_complete_repetitions}{[n]} complete \resultplural{v4x.x6.minimum_complete_repetitions}{repetition}{repetitions} against D1 (\texttt{no reading}).
Counts and Wilson intervals, where available, are retained without a coverage verdict.

\paragraph{X6: reference precision.}
Reference precision left \resultnum{v4x.x6.x6a.count.inconclusive-truth-precision}{[n]} X6a \resultplural{v4x.x6.x6a.count.inconclusive-truth-precision}{reference}{references} and \resultnum{v4x.x6.x6b.count.inconclusive-truth-precision}{[n]} X6b \resultplural{v4x.x6.x6b.count.inconclusive-truth-precision}{reference}{references} inconclusive against D1 (\texttt{inconclusive-truth-precision}).
Available sensitivity results and the original nominal reading are retained (Table~\ref{tab:v4x-x6-precision}).

% X6-UNAVAILABLE: select when x6a.band.vs_d1 == null, or per cell with vs_d1 null.
\paragraph{X6: unavailable procedural truth.}
Where the procedural truth was unavailable, no D1 coverage reference was calculated.
The truth estimates, errors, gaps and operational inventory remain recorded without a substituted truth or coverage denominator.

\paragraph{X6: conservative flag.}
A Wilson lower endpoint above \resultnum{v4x.protocol.x6.conservative}{[value]} produced the conservative flag in \resultnum{v4x.x6.x6a.count.conservative}{[n]} X6a coverage \resultplural{v4x.x6.x6a.count.conservative}{reference}{references} and \resultnum{v4x.x6.x6b.count.conservative}{[n]} D1 coverage \resultplural{v4x.x6.x6b.count.conservative}{reference}{references}.
The flag applies only to each stated denominator.

% X6-TRUTH: always; no coverage verdict if the D1 truth is unavailable against D1.
\paragraph{X6: procedural and approximate population truth.}
The primary truth conditions on successful D1 pipelines and includes matching with \resultnum{v4x.x3.sizes.matching_prose}{[n]} \resultplural{v4x.x3.sizes.matching_prose}{pair}{pairs} and the declared \resultnum{v4x.x3.sizes.retry_matching_prose}{[n]}-pair retries.
The secondary truth uses approximate population matching with \resultnum{v4x.x6.sizes.population_pairs_prose}{[n]} \resultplural{v4x.x6.sizes.population_pairs_prose}{pair}{pairs}.
Table~\ref{tab:v4x-x6-truth} reports retry fractions and truth Monte Carlo errors, against the declared precision target of \resultnum{v4x.x6.truth_se_target}{[value]}.
The secondary truth's Monte Carlo error omits the population-fit approximation error.
The calibration exposure at the declared X6a signal share remains relevant to adaptation even though this analysis does not read boundaries.

\paragraph{X6: matching diagnostics.}
Table~\ref{tab:v4x-x6-diagnostics} reports standardized errors and their recorded bias flags, the D5 mean shift, empirical matching bias and the procedural-minus-population truth difference. It also reports fit residuals, amplification and the ratio of mean squared error to mean squared standard error.
These diagnostics retain their own success conditioning and reference-precision limits.

% X7-SCOPE: always with X7; the comparison and its common limits, stated once.
\paragraph{X7: scope.}
X7 compares an LDA on expanded features with the frozen-feature LDA on paired exploratory seeds at every response signal share (Tables~\ref{tab:v4x-x7}, \ref{tab:v4x-x7-boundaries} and~\ref{tab:v4x-x7-deltas}); errors and gaps of the expanded-feature arm are recorded separately from those of the frozen-feature arm.
Changing feature dimension changes estimation and regularization as well as information, so a richer fitted LDA establishes neither an information ceiling nor Bayes optimality.

% X7: one branch per positive unit count, retaining invalid/censored states.
\paragraph{X7: gain under the paired criterion.}
In \resultnum{v4x.x7.count.far}{[n]} \resultplural{v4x.x7.count.far}{unit}{units}, the expanded-feature LDA gained relative to the frozen-feature LDA under the declared paired criterion (\texttt{far}).
In these units, either the expanded boundary interval's upper endpoint was below the frozen-feature point boundary, or the paired AUC-difference intervals were positive at every grid strength from the expanded first boundary up to, but excluding, the frozen-feature first boundary (Tables~\ref{tab:v4x-x7-boundaries} and \ref{tab:v4x-x7-deltas}).
The exact applicable predicate and any censoring are retained by unit.
For this predicate, a frozen-feature first boundary not attained without point censoring is treated as beyond the grid.

\paragraph{X7: local equivalence near the boundary.}
In \resultnum{v4x.x7.count.close}{[n]} \resultplural{v4x.x7.count.close}{unit}{units}, the two LDAs showed local equivalence within $\pm\resultnum{v4x.x7.bootstrap.margin}{[value]}$ near their common grid boundary (\texttt{close}).
The first point boundaries coincided and the paired intervals at that strength and its grid neighbours lay inside the declared margin.

\paragraph{X7: inconclusive comparison.}
The declared paired comparison was inconclusive in \resultnum{v4x.x7.count.inconclusive}{[n]} \resultplural{v4x.x7.count.inconclusive}{unit}{units} (\texttt{inconclusive}); the two boundaries and paired contrasts remain tabulated.

\paragraph{X7: invalid comparison.}
There \resultplural{v4x.x7.count.invalid}{was}{were} \resultnum{v4x.x7.count.invalid}{[n]} invalid \resultplural{v4x.x7.count.invalid}{unit}{units} (\texttt{invalid}) because a required zero-strength gate was invalid.
\resultplural{v4x.x7.count.invalid}{Its}{Their} status remains beside any recorded curves, without a valid boundary interpretation.

\paragraph{X7: censored comparison.}
There \resultplural{v4x.x7.count.censored}{was}{were} \resultnum{v4x.x7.count.censored}{[n]} censored unit \resultplural{v4x.x7.count.censored}{reading}{readings} (\texttt{censored}).
Point and interval censoring are retained separately for first and sustained boundaries; a missing crossing required by the paired rule is not converted into an AUC below target.

\paragraph{X7: unavailable comparison.}
There \resultplural{v4x.x7.count.unavailable}{was}{were} \resultnum{v4x.x7.count.unavailable}{[n]} unavailable \resultplural{v4x.x7.count.unavailable}{reading}{readings} (\texttt{unavailable}). \resultplural{v4x.x7.count.unavailable}{Its frozen-arm}{Their frozen-arm} and expanded-arm errors remain recorded.

\paragraph{X7: incomplete comparison.}
There \resultplural{v4x.x7.count.incomplete}{was}{were} \resultnum{v4x.x7.count.incomplete}{[n]} incomplete \resultplural{v4x.x7.count.incomplete}{unit}{units} at the operational cut (\texttt{incomplete}).
The cell inventory recorded \resultnum{v4x.x7.inventory.recorded}{[n]} \resultplural{v4x.x7.inventory.recorded}{cell}{cells}, with \resultnum{v4x.x7.inventory.incomplete}{[n]} incomplete \resultplural{v4x.x7.inventory.incomplete}{cell}{cells}, \resultnum{v4x.x7.inventory.errors}{[n]} frozen-arm \resultplural{v4x.x7.inventory.errors}{error}{errors} and \resultnum{v4x.x7.inventory.gaps}{[n]} frozen-arm fit \resultplural{v4x.x7.inventory.gaps}{gap}{gaps}.

% POST-HEADING: only with at least one post-release reading below.
\subsection*{Exploratory readings beside the released boundary}

% POST-AVAILABLE: each component has its own availability count.
\paragraph{Available local readings.}
After release, the first confirmatory point boundary selected an X3 cell with an available contrast at \resultnum{v4x.comparisons.count.x3.available}{[n]} response signal \resultplural{v4x.comparisons.count.x3.available}{share}{shares}. Both conditional X4 references were available at \resultnum{v4x.comparisons.count.x4.available}{[n]} \resultplural{v4x.comparisons.count.x4.available}{share}{shares} (Tables~\ref{tab:v4x-comparisons} and \ref{tab:v4x-selected-comparisons}).
Selection used the first point boundary and did not interpolate strengths or select by the sustained boundary. A missing or gap-censored first point selects no cell, and a finite first point with a censored interval endpoint selects its exact cell with that limitation retained.
The original X3 label, paired interval, X4 denominators, size flags and quantity labels remain subject to the same local reading limits.

% POST-PARTIAL/UNAVAILABLE: retain original reasons, never treat absence as zero.
\paragraph{Partial and unavailable readings.}
For X4, \resultnum{v4x.comparisons.count.x4.partial}{[n]} response signal \resultplural{v4x.comparisons.count.x4.partial}{share had}{shares had} a partial reading, and \resultnum{v4x.comparisons.count.x4.unavailable}{[n]} \resultplural{v4x.comparisons.count.x4.unavailable}{had an unavailable reading}{had unavailable readings}. For X3, \resultnum{v4x.comparisons.count.x3.unavailable}{[n]} \resultplural{v4x.comparisons.count.x3.unavailable}{share had an unavailable reading}{shares had unavailable readings} (Table~\ref{tab:v4x-comparison-states}).
Reasons distinguish a missing or censored confirmatory first boundary, a zero-strength boundary outside the X3 declaration, for which no neighbouring positive-strength cell is substituted, a missing aggregate or cell, a fit gap and a recorded error. Outside-domain states are recorded separately and are excluded from the unavailable X3 counts.
Local availability retains the global main-curve validity mark.

% POST-X1B/X7: side by side only; X7 is joined by r from its own aggregate.
\paragraph{Twin and expanded-feature boundaries.}\phantomsection\label{supp:x1b-release}
In the supplied post-release aggregate, the twin first boundary was available at \resultnum{v4x.comparisons.count.x1b.available}{[n]} response signal \resultplural{v4x.comparisons.count.x1b.available}{share}{shares}; at the other \resultnum{v4x.comparisons.count.x1b.unavailable}{[n]} it was censored at a gap, not reached on the complete eligible grid, or undefined for an invalid or unavailable unit (Tables~\ref{tab:v4x-comparison-states} and~\ref{tab:v4x-boundary-comparisons}). The full supplied X1b display retains every declared signal share, first and sustained records, intervals, statuses and separate censoring; an invalid or unavailable record is not replaced by the archived numerical account in the fixed supplement.
In the supplied post-release aggregate, the twin first boundary was available at \resultnum{v4x.comparisons.count.x1b.available}{[n]} response signal \resultplural{v4x.comparisons.count.x1b.available}{share}{shares} (Tables~\ref{tab:v4x-comparison-states} and~\ref{tab:v4x-boundary-comparisons}). The full supplied X1b display retains every declared signal share, first and sustained records, intervals, statuses and separate censoring; an invalid or unavailable record is not replaced by the archived numerical account in the fixed supplement.
In the supplied post-release aggregate, no twin first boundary was available. At all \resultnum{v4x.comparisons.count.x1b.unavailable}{[n]} response signal \resultplural{v4x.comparisons.count.x1b.unavailable}{share}{shares}, it was censored at a gap, not reached on the complete eligible grid, or undefined for an invalid or unavailable unit (Tables~\ref{tab:v4x-comparison-states} and~\ref{tab:v4x-boundary-comparisons}). The full supplied X1b display retains every declared signal share, first and sustained records, intervals, statuses and separate censoring; an invalid or unavailable record is not replaced by the archived numerical account in the fixed supplement.
Table~\ref{tab:v4x-boundary-comparisons} places the twin and both exploratory LDA first boundaries beside the released $a^*(r)$, retaining their own intervals, statuses and censoring in the accompanying tables and catalog.
Table~\ref{tab:v4x-stress-comparisons} records the misspecification comparisons; Table~\ref{tab:v4x-x1b-boundaries} retains the full twin boundary records. These displays add no between-curve test across the different seed sets and do not isolate network dependence.

% POST-X5: select the applicable branch per condition and r from the table.
\paragraph{Misspecification equivalent strength: available.}
For \resultnum{v4x.comparisons.count.x5.equivalent.available}{[n]} available X5 condition--signal-share \resultplural{v4x.comparisons.count.x5.equivalent.available}{reading}{readings}, Table~\ref{tab:v4x-stress-comparisons} reports the first declared grid strength whose confirmatory point AUC equals or exceeds the misspecification AUC.

\paragraph{Misspecification equivalent strength: censored.}
A gap before the equivalent-strength attainment left \resultnum{v4x.comparisons.count.x5.equivalent.censored}{[n]} condition--signal-share \resultplural{v4x.comparisons.count.x5.equivalent.censored}{reading}{readings} censored at the recorded strength.

\paragraph{Misspecification equivalent strength: above the grid.}
There was no attainment on a complete eligible grid for \resultnum{v4x.comparisons.count.x5.equivalent.above-grid}{[n]} \resultplural{v4x.comparisons.count.x5.equivalent.above-grid}{reading}{readings}, retained as above the grid without interpolation or extrapolation.

\paragraph{Misspecification interval: overlap.}
At the released first boundary, \resultnum{v4x.comparisons.count.x5.interval.overlap}{[n]} misspecification \resultplural{v4x.comparisons.count.x5.interval.overlap}{interval}{intervals} overlapped the confirmatory band; touching an endpoint counts as overlap.

\paragraph{Misspecification interval: above.}
At the released first boundary, \resultnum{v4x.comparisons.count.x5.interval.above}{[n]} misspecification \resultplural{v4x.comparisons.count.x5.interval.above}{interval}{intervals} lay above the confirmatory band.

\paragraph{Misspecification interval: below.}
At the released first boundary, \resultnum{v4x.comparisons.count.x5.interval.below}{[n]} misspecification \resultplural{v4x.comparisons.count.x5.interval.below}{interval}{intervals} lay below the confirmatory band.

\paragraph{Misspecification comparison: unavailable.}
An invalid or incomplete main unit, a missing or censored first boundary, or an unavailable X5 estimate retains an unavailable reading and its original reason.
The counts were \resultnum{v4x.comparisons.count.x5.equivalent.unavailable}{[n]} unavailable equivalent \resultplural{v4x.comparisons.count.x5.equivalent.unavailable}{strength}{strengths} and \resultnum{v4x.comparisons.count.x5.interval.unavailable}{[n]} unavailable interval \resultplural{v4x.comparisons.count.x5.interval.unavailable}{comparison}{comparisons}.
An available equivalent strength can coexist with an unavailable interval comparison when the confirmatory first boundary is not attained.

% TABLE INPUTS: fixed filenames, emitted by the export. No hand-written values.
\input{generated/exploratory-status}
\input{generated/exploratory-x3}
\input{generated/exploratory-x3-attempts}
\input{generated/exploratory-x4a}
\input{generated/exploratory-x4b}
\input{generated/exploratory-x4-normal}
\input{generated/exploratory-x4-calibration}
\input{generated/exploratory-x4-failures}
\input{generated/exploratory-x4-paired}
\input{generated/exploratory-x6-inventory}
\input{generated/exploratory-x6}
\input{generated/exploratory-x6-denominators}
\input{generated/exploratory-x6-truth}
\input{generated/exploratory-x6-precision}
\input{generated/exploratory-x6-diagnostics}
\input{generated/exploratory-x7}
% X7 boundaries are included once by main.tex; retain that full display.
\input{generated/exploratory-x7-deltas}
\input{generated/exploratory-comparisons}
\input{generated/exploratory-selected-comparisons}
\input{generated/exploratory-comparison-states}
\input{generated/exploratory-boundary-comparisons}
\input{generated/exploratory-stress-comparisons}
